% Options for packages loaded elsewhere
\PassOptionsToPackage{unicode}{hyperref}
\PassOptionsToPackage{hyphens}{url}
\PassOptionsToPackage{space}{xeCJK}
\documentclass[
]{article}
\usepackage{xcolor}
\usepackage{amsmath,amssymb}
\setcounter{secnumdepth}{-\maxdimen} % remove section numbering
\usepackage{iftex}
\ifPDFTeX
  \usepackage[T1]{fontenc}
  \usepackage[utf8]{inputenc}
  \usepackage{textcomp} % provide euro and other symbols
\else % if luatex or xetex
  \usepackage{unicode-math} % this also loads fontspec
  \defaultfontfeatures{Scale=MatchLowercase}
  \defaultfontfeatures[\rmfamily]{Ligatures=TeX,Scale=1}
\fi
\usepackage{lmodern}
\ifPDFTeX\else
  % xetex/luatex font selection
  \setmainfont[]{Times New Roman}
  \setsansfont[]{SimHei}
  \setmonofont[]{SimKai}
  \ifXeTeX
    \usepackage{xeCJK}
    \setCJKmainfont[]{SimSun}
  \fi
  \ifLuaTeX
    \usepackage[]{luatexja-fontspec}
    \setmainjfont[]{SimSun}
  \fi
\fi
% Use upquote if available, for straight quotes in verbatim environments
\IfFileExists{upquote.sty}{\usepackage{upquote}}{}
\IfFileExists{microtype.sty}{% use microtype if available
  \usepackage[]{microtype}
  \UseMicrotypeSet[protrusion]{basicmath} % disable protrusion for tt fonts
}{}
\makeatletter
\@ifundefined{KOMAClassName}{% if non-KOMA class
  \IfFileExists{parskip.sty}{%
    \usepackage{parskip}
  }{% else
    \setlength{\parindent}{0pt}
    \setlength{\parskip}{6pt plus 2pt minus 1pt}}
}{% if KOMA class
  \KOMAoptions{parskip=half}}
\makeatother
\usepackage{longtable,booktabs,array}
\newcounter{none} % for unnumbered tables
\usepackage{multirow}
\usepackage{calc} % for calculating minipage widths
% Correct order of tables after \paragraph or \subparagraph
\usepackage{etoolbox}
\makeatletter
\patchcmd\longtable{\par}{\if@noskipsec\mbox{}\fi\par}{}{}
\makeatother
% Allow footnotes in longtable head/foot
\IfFileExists{footnotehyper.sty}{\usepackage{footnotehyper}}{\usepackage{footnote}}
\makesavenoteenv{longtable}
\ifLuaTeX
  \usepackage{luacolor}
  \usepackage[soul]{lua-ul}
\else
  \usepackage{soul}
  \ifXeTeX
    % soul's \st doesn't work for CJK:
    \usepackage{xeCJKfntef}
    \renewcommand{\st}[1]{\sout{#1}}
  \fi
\fi
\setlength{\emergencystretch}{3em} % prevent overfull lines
\providecommand{\tightlist}{%
  \setlength{\itemsep}{0pt}\setlength{\parskip}{0pt}}
\usepackage{bookmark}
\IfFileExists{xurl.sty}{\usepackage{xurl}}{} % add URL line breaks if available
\urlstyle{same}
\hypersetup{
  hidelinks,
  pdfcreator={LaTeX via pandoc}}

\author{}
\date{}

\begin{document}

\textbf{计算机学院本科毕业论文（设计）开题答辩评分表}

{\def\LTcaptype{none} % do not increment counter
\begin{longtable}[]{@{}
  >{\centering\arraybackslash}p{(\linewidth - 12\tabcolsep) * \real{0.1037}}
  >{\centering\arraybackslash}p{(\linewidth - 12\tabcolsep) * \real{0.1338}}
  >{\centering\arraybackslash}p{(\linewidth - 12\tabcolsep) * \real{0.1637}}
  >{\centering\arraybackslash}p{(\linewidth - 12\tabcolsep) * \real{0.1487}}
  >{\centering\arraybackslash}p{(\linewidth - 12\tabcolsep) * \real{0.1785}}
  >{\centering\arraybackslash}p{(\linewidth - 12\tabcolsep) * \real{0.1636}}
  >{\centering\arraybackslash}p{(\linewidth - 12\tabcolsep) * \real{0.1082}}@{}}
\toprule\noalign{}
\begin{minipage}[b]{\linewidth}\centering
\textbf{姓名}
\end{minipage} & \begin{minipage}[b]{\linewidth}\centering
\end{minipage} & \begin{minipage}[b]{\linewidth}\centering
\textbf{学号}
\end{minipage} & \begin{minipage}[b]{\linewidth}\centering
\end{minipage} & \begin{minipage}[b]{\linewidth}\centering
\textbf{年级}
\end{minipage} &
\multicolumn{2}{>{\centering\arraybackslash}p{(\linewidth - 12\tabcolsep) * \real{0.2717} + 2\tabcolsep}@{}}{%
\begin{minipage}[b]{\linewidth}\centering
\end{minipage}} \\
\midrule\noalign{}
\endhead
\bottomrule\noalign{}
\endlastfoot
\textbf{专业} & & \textbf{班级} & & \textbf{指导教师} &
\multicolumn{2}{>{\centering\arraybackslash}p{(\linewidth - 12\tabcolsep) * \real{0.2717} + 2\tabcolsep}@{}}{%
} \\
\multicolumn{2}{@{}>{\centering\arraybackslash}p{(\linewidth - 12\tabcolsep) * \real{0.2374} + 2\tabcolsep}}{%
论文(设计)题目} &
\multicolumn{5}{>{\centering\arraybackslash}p{(\linewidth - 12\tabcolsep) * \real{0.7626} + 8\tabcolsep}@{}}{%
} \\
\multicolumn{2}{@{}>{\centering\arraybackslash}p{(\linewidth - 12\tabcolsep) * \real{0.2374} + 2\tabcolsep}}{%
\textbf{评价内容}} &
\multicolumn{4}{>{\centering\arraybackslash}p{(\linewidth - 12\tabcolsep) * \real{0.6544} + 6\tabcolsep}}{%
\textbf{评价标准}} & \textbf{得分} \\
\multicolumn{2}{@{}>{\centering\arraybackslash}p{(\linewidth - 12\tabcolsep) * \real{0.2374} + 2\tabcolsep}}{%
\multirow{4}{=}{\centering\arraybackslash \textbf{设计方案}

\textbf{（30分）}}} &
\multicolumn{4}{>{\centering\arraybackslash}p{(\linewidth - 12\tabcolsep) * \real{0.6544} + 6\tabcolsep}}{%
优秀（30-27分）：能提出最优的技术方案，考虑安全、健康、法律、文化以及环境等因素，可行性分析论证充分，论证过程完善。}
& \multirow{4}{=}{\centering\arraybackslash } \\
& &
\multicolumn{4}{>{\centering\arraybackslash}p{(\linewidth - 12\tabcolsep) * \real{0.6544} + 6\tabcolsep}}{%
良好（26-23分）：能提出合理的技术方案，考虑安全、健康、法律、文化以及环境等因素，可行性分析论证充分，过程完整。} \\
& &
\multicolumn{4}{>{\centering\arraybackslash}p{(\linewidth - 12\tabcolsep) * \real{0.6544} + 6\tabcolsep}}{%
及格（22-18）：能提出合理的技术方案，考虑安全、健康、法律、文化以及环境等因素，进行了可行性分析论证。} \\
& &
\multicolumn{4}{>{\centering\arraybackslash}p{(\linewidth - 12\tabcolsep) * \real{0.6544} + 6\tabcolsep}}{%
不及格（17-0）：无技术方案的合理性和可行性论证。} \\
\multicolumn{2}{@{}>{\centering\arraybackslash}p{(\linewidth - 12\tabcolsep) * \real{0.2374} + 2\tabcolsep}}{%
\multirow{4}{=}{\centering\arraybackslash \textbf{沟通表达}

\textbf{（30分）}}} &
\multicolumn{4}{>{\centering\arraybackslash}p{(\linewidth - 12\tabcolsep) * \real{0.6544} + 6\tabcolsep}}{%
优秀（30-27分）：开题报告和文献综述观点鲜明、条理清晰、语言准确、格式规范，答辩陈述清晰，回答问题对答如流。}
& \multirow{4}{=}{\centering\arraybackslash } \\
& &
\multicolumn{4}{>{\centering\arraybackslash}p{(\linewidth - 12\tabcolsep) * \real{0.6544} + 6\tabcolsep}}{%
良好（26-23）：开题报告和文献综述层次清晰、文字通顺、格式正确，答辩陈述清晰，回答问题准确。} \\
& &
\multicolumn{4}{>{\centering\arraybackslash}p{(\linewidth - 12\tabcolsep) * \real{0.6544} + 6\tabcolsep}}{%
及格（22-18）：开题报告和文献综述结构合理、文字通顺、格式错误少，答辩简洁，能够回答问题。} \\
& &
\multicolumn{4}{>{\centering\arraybackslash}p{(\linewidth - 12\tabcolsep) * \real{0.6544} + 6\tabcolsep}}{%
不及格（17-0）：开题报告和文献综述逻辑不清，格式错误多，陈述不清，不能回答提问。} \\
\multicolumn{2}{@{}>{\centering\arraybackslash}p{(\linewidth - 12\tabcolsep) * \real{0.2374} + 2\tabcolsep}}{%
\multirow{4}{=}{\centering\arraybackslash \textbf{自主学习}

\textbf{（40分）}}} &
\multicolumn{4}{>{\centering\arraybackslash}p{(\linewidth - 12\tabcolsep) * \real{0.6544} + 6\tabcolsep}}{%
优秀（40-36分）：选题涉及技术前沿，自主学习了大量前沿知识和技术，设计方案中很好的运用了新知识新技术。}
& \multirow{4}{=}{\centering\arraybackslash } \\
& &
\multicolumn{4}{>{\centering\arraybackslash}p{(\linewidth - 12\tabcolsep) * \real{0.6544} + 6\tabcolsep}}{%
良好（35-30）：选题较新颖，结合选题学习了一些新知识新技术，设计方案运用了新知识新技术。} \\
& &
\multicolumn{4}{>{\centering\arraybackslash}p{(\linewidth - 12\tabcolsep) * \real{0.6544} + 6\tabcolsep}}{%
及格（29-24）：选题有新意，结合选题学习了新知识新技术，设计方案中运用了新知识新技术。} \\
& &
\multicolumn{4}{>{\centering\arraybackslash}p{(\linewidth - 12\tabcolsep) * \real{0.6544} + 6\tabcolsep}}{%
不及格（23-0）：选题新意不足，自学能力一般，设计方案较少运用新知识技术。} \\
\multicolumn{6}{@{}>{\centering\arraybackslash}p{(\linewidth - 12\tabcolsep) * \real{0.8918} + 10\tabcolsep}}{%
\textbf{综合得分}} & \\
\multicolumn{2}{@{}>{\centering\arraybackslash}p{(\linewidth - 12\tabcolsep) * \real{0.2374} + 2\tabcolsep}}{%
\textbf{开题答辩小组意见}} &
\multicolumn{5}{>{\centering\arraybackslash}p{(\linewidth - 12\tabcolsep) * \real{0.7626} + 8\tabcolsep}@{}}{%
\textbf{\hl{（评语）：}}

评审小组教师签字：\hl{小组成员签名}} \\
\end{longtable}
}

\end{document}
